%-----------------------------------------------------------------------------------------------------------------------------------------------%
%	The MIT License (MIT)
%
%	Copyright (c) 2015 Jan Küster
%
%	Permission is hereby granted, free of charge, to any person obtaining a copy
%	of this software and associated documentation files (the "Software"), to deal
%	in the Software without restriction, including without limitation the rights
%	to use, copy, modify, merge, publish, distribute, sublicense, and/or sell
%	copies of the Software, and to permit persons to whom the Software is
%	furnished to do so, subject to the following conditions:
%	
%	THE SOFTWARE IS PROVIDED "AS IS", WITHOUT WARRANTY OF ANY KIND, EXPRESS OR
%	IMPLIED, INCLUDING BUT NOT LIMITED TO THE WARRANTIES OF MERCHANTABILITY,
%	FITNESS FOR A PARTICULAR PURPOSE AND NONINFRINGEMENT. IN NO EVENT SHALL THE
%	AUTHORS OR COPYRIGHT HOLDERS BE LIABLE FOR ANY CLAIM, DAMAGES OR OTHER
%	LIABILITY, WHETHER IN AN ACTION OF CONTRACT, TORT OR OTHERWISE, ARISING FROM,
%	OUT OF OR IN CONNECTION WITH THE SOFTWARE OR THE USE OR OTHER DEALINGS IN
%	THE SOFTWARE.
%	
%
%-----------------------------------------------------------------------------------------------------------------------------------------------%


%============================================================================%
%
%	DOCUMENT DEFINITION
%
%============================================================================%

%we use article class because we want to fully customize the page and dont use a cv template
\documentclass[10pt,A4]{article}	


%----------------------------------------------------------------------------------------
%	ENCODING
%----------------------------------------------------------------------------------------

%we use utf8 since we want to build from any machine
\usepackage[utf8]{inputenc}		

%----------------------------------------------------------------------------------------
%	LOGIC
%----------------------------------------------------------------------------------------

% provides \isempty test
\usepackage{xifthen}

%----------------------------------------------------------------------------------------
%	FONT
%----------------------------------------------------------------------------------------

% some tex-live fonts - choose your own

%\usepackage[defaultsans]{droidsans}
%\usepackage[default]{comfortaa}
%\usepackage{cmbright}
\usepackage[default]{raleway}
%\usepackage{fetamont}
%\usepackage[default]{gillius}
%\usepackage[light,math]{iwona}
%\usepackage[thin]{roboto} 

% set font default
\renewcommand*\familydefault{\sfdefault} 	
\usepackage[T1]{fontenc}

% more font size definitions
\usepackage{moresize}		


%----------------------------------------------------------------------------------------
%	PAGE LAYOUT  DEFINITIONS
%----------------------------------------------------------------------------------------

%debug page outer frames
%\usepackage{showframe}			


%define page styles using geometry
\usepackage[a4paper]{geometry}		

% for example, change the margins to 2 inches all round
\geometry{top=1.25cm, bottom=-.6cm, left=1.5cm, right=1.5cm} 	

%use customized header
\usepackage{fancyhdr}				
\pagestyle{fancy}

%less space between header and content
\setlength{\headheight}{-5pt}		


%customize entries left, center and right
%\lhead{}
%\chead{}
%\rhead{}


%indentation is zero
\setlength{\parindent}{0mm}

%----------------------------------------------------------------------------------------
%	TABLE /ARRAY DEFINITIONS
%---------------------------------------------------------------------------------------- 

%for layouting tables
\usepackage{multicol}			
\usepackage{multirow}

%extended aligning of tabular cells
\usepackage{array}

\newcolumntype{x}[1]{%
>{\raggedleft\hspace{0pt}}p{#1}}%


%----------------------------------------------------------------------------------------
%	GRAPHICS DEFINITIONS
%---------------------------------------------------------------------------------------- 

%for header image
\usepackage{graphicx}

%for floating figures
\usepackage{wrapfig}
\usepackage{float}
%\floatstyle{boxed} 
%\restylefloat{figure}

%for drawing graphics		
\usepackage{tikz}				
\usetikzlibrary{shapes, backgrounds,mindmap, trees}


%----------------------------------------------------------------------------------------
%	Color DEFINITIONS
%---------------------------------------------------------------------------------------- 

\usepackage{color}

%accent color
\definecolor{sectcol}{RGB}{0,150,255}

%dark background color
\definecolor{bgcol}{RGB}{110,110,110}

%light background / accent color
\definecolor{softcol}{RGB}{225,225,225}


%============================================================================%
%
%
%	DEFINITIONS
%
%
%============================================================================%

%----------------------------------------------------------------------------------------
% 	HEADER
%----------------------------------------------------------------------------------------

% remove top header line
\renewcommand{\headrulewidth}{0pt} 

%remove botttom header line
\renewcommand{\footrulewidth}{0pt}	  	

%remove pagenum
\renewcommand{\thepage}{}	

%remove section num		
\renewcommand{\thesection}{}			

%----------------------------------------------------------------------------------------
% 	ARROW GRAPHICS in Tikz
%----------------------------------------------------------------------------------------

% a six pointed arrow poiting to the left
\newcommand{\tzlarrow}{(0,0) -- (0.2,0) -- (0.3,0.2) -- (0.2,0.4) -- (0,0.4) -- (0.1,0.2) -- cycle;}	

% include the left arrow into a tikz picture
% param1: fill color
%
\newcommand{\larrow}[1]
{\begin{tikzpicture}[scale=0.58]
	 \filldraw[fill=#1!100,draw=#1!100!black]  \tzlarrow
 \end{tikzpicture}
}

% a six pointed arrow poiting to the right
\newcommand{\tzrarrow}{ (0,0.2) -- (0.1,0) -- (0.3,0) -- (0.2,0.2) -- (0.3,0.4) -- (0.1,0.4) -- cycle;}

% include the right arrow into a tikz picture
% param1: fill color
%
\newcommand{\rarrow}
{\begin{tikzpicture}[scale=0.7]
	\filldraw[fill=sectcol!100,draw=sectcol!100!black] \tzrarrow
 \end{tikzpicture}
}



%----------------------------------------------------------------------------------------
%	custom sections
%----------------------------------------------------------------------------------------

% create a coloured box with arrow and title as cv section headline
% param 1: section title
%
\newcommand{\cvsection}[1]
{
	\begin{center}
		\large\textcolor{sectcol}{\textbf{#1}}
	\end{center}
}

%create a coloured arrow with title as cv meta section section
% param 1: meta section title
%
\newcommand{\metasection}[2]
{
%\begin{tabular*}{1\textwidth}{r r}
\footnotesize{#2} \hspace*{\fill} \footnotesize{#1}\\[1pt]
%\end{tabular*}
}

%----------------------------------------------------------------------------------------
%	 CV EVENT
%----------------------------------------------------------------------------------------

% creates a stretched box as cv entry headline followed by some paragraphs about 
% the work you did
% param 1:	event time i.e. 2014 or 2011-2014 etc.
% param 2:	event name (what did you do?)
% param 3:	institution (where did you work / study)
% param 4:	list of paragraphs
%
\newcommand{\cvevent}[4]
{

\begin{tabular*}{1\textwidth}{p{13.6cm}  x{3.9cm}}
	\textbf{#2} - \textcolor{bgcol}{#3} &   \vspace{2.5pt}\textcolor{sectcol}{#1}
\end{tabular*}

\vspace{-8pt}
\textcolor{softcol}{\hrule}
\vspace{6pt}

	\foreach \desc in {#4}{
		$\cdot$ \desc\\[3pt]
	}
	
\vspace{3pt}
}

% creates a stretched box as 
\newcommand{\cveventmeta}[2]
{
	\mbox{\mystrut \hspace{87pt}\textit{#1}}\\
	#2
}

%----------------------------------------------------------------------------------------
% CUSTOM STRUT FOR EMPTY BOXES
%----------------------------------------- -----------------------------------------------
\newcommand{\mystrut}{\rule[-.3\baselineskip]{0pt}{\baselineskip}}

%----------------------------------------------------------------------------------------
% CUSTOM LOREM IPSUM
%----------------------------------------------------------------------------------------
\newcommand{\lorem}
{Lorem ipsum dolor sit amet, consectetur adipiscing elit. Donec a diam lectus.}



%============================================================================%
%
%
%
%	DOCUMENT CONTENT
%
%
%
%============================================================================%
\begin{document}


%use our custom fancy header definitions
\pagestyle{fancy}	





%----------------------------------------------------------------------------------------
%	HEADER IMAGE
%----------------------------------------------------------------------------------------

%\begin{figure}[H]
%\begin{flushright}
%	\includegraphics[trim= 320 130 460 210,clip,width=0.2\linewidth]{myfoto.jpg}	%trimming relative to image size!
%\end{flushright}
%\end{figure}


%---------------------------------------------------------------------------------------
%	TITLE HEADLINE
%----------------------------------------------------------------------------------------
\vspace{-8pt}
\begin{center}
	\HUGE \textsc{Yuna Son} \textcolor{sectcol}{\rule[-1mm]{1mm}{0.9cm}} \textsc{Resume}\\[2pt]
	\Ph.D. Candidate in Biomedical Sciences
\end{center}



\vspace{6pt}


%---------------------------------------------------------------------------------------
%	META SECTION
%----------------------------------------------------------------------------------------
\metasection{Tallahassee, Florida, USA}{\textbf{Status:} Ph.D. candidate in Florida State University}
\metasection{yuna.son@med.fsu.edu}{\textbf{Research:} Fibrosis, Keloids, Translational research, Genomics} 
\metasection{github.com/sunabcun}{\textbf{Techs:} R, Python, Bash, Linux for HPC}
\metasection{850-559-9118}{\textbf{Skills:} RNAseq, WGS, WGBS, Histology, Molecular and Cellular experiments, RNA/DNA works}
\vspace{-2pt}
\textcolor{softcol}{\hrule}
\vspace{6pt}

\normalsize

%---------------------------------------------------------------------------------------
%	SUMMARAY (optional)
%----------------------------------------------------------------------------------------
\vspace{-6pt}
\cvsection{Summary}
More than 7 years of academic lab experiences; Trained and skilled in molecular biology, histology, and RNA/DNA works; 
2+ years of NGS library preparation such as RNAseq, WGS, and WGBS; 2+ years of experience in analyzing large genomic data set using statistical/computational approaches;
Strong problem solving and data analysis skills.\\


%============================================================================%
%
%	CV SECTIONS AND EVENTS (MAIN CONTENT)
%
%============================================================================%

%---------------------------------------------------------------------------------------
%	EXPERIENCE
%----------------------------------------------------------------------------------------
\cvsection{Research Experience}

\cvevent{2016 - present}{Research Assistant}{Florida State University}{
Unveiling the molecular mechanisms involved in fibrotic diseases using the fibrotic skin lesions known as keloids.
	{Evaluated various fractionation and dosing schemes of radiation to prevent keloid recurrence in vitro and patient data. (DOI: 10.1016/j.jaad.2020.06.023.)},
	{Provided further understanding of how radiation can reduce keloids proliferation using histology and immunofluorescence microscopy.},
	{Developed methods to predict steroid responsiveness in non-invasive approaches. (Application-62/991,423)},
	{Determined genetic and epigenetic mechanisms involved in the formation and therapeutic responses of keloids using genome wide sequencing based approaches.}
}

%
\cvevent{2014 - 2016}{Research Assistant}{Korea University}{
Determining the therapeutic effects of drug candidates for fibrosis using mouse models.
	{Generated bleomycin-induced lung fibrosis mouse models.},
	{Routinely performed IV injections to deliver a drug candidate.},
	{Determined the clinically beneficial effects of the drug candidate by histology and molecular analysis.}
}

%\textcolor{softcol}{\hrule}

%

%---------------------------------------------------------------------------------------
%	Education SECTION
%--------------------------------------------------------------------------------------
\cvsection{Education}

\cvevent{2016-Present}{Ph.D. Candidate in Biomedical Sciences}{Florida State University}{
	{Supervisor: Akash Gunjan, Ph.D.}}

%\textcolor{softcol}{\hrule}

%
\cvevent{2016}{M.S. in Biomedical Sciences}{Korea University}{
	{Master thesis: Fusion protein of retinol-binding protein and albumin domain III reduces bleomycin-induced lung fibrosis}
}

%\textcolor{softcol}{\hrule}

%
\cvevent{2014}{B.S. in Biotechnology and Life Sciences}{Kyungpook National University}


%---------------------------------------------------------------------------------------
%	Publications SECTION
%--------------------------------------------------------------------------------------
\cvsection{Publications}{
{\textbf{Son, Y.}, Phillips, E., Price, K., Rosenberg, L., Stefanovic, B., Wolfe, C., Shaath, T., Om,
A., Cohen, G., Gunjan, A. Treatment of keloids with a single dose of low energy superficial X-ray radiation to prevent recurrence after surgical excision: an in vitro and in vivo study, Journal of the American Academy of Dermatology, 2020, DOI: 10.1016/j.jaad.2020.06.023. (Impact factor: 8.277)},
{Hammack, C., Ogden, S. C., Madden, J., Medina, A., Xu, C., Phillips, E., \textbf{Son, Y.},
Sherman, A., Giovinazzi, S., Didier, R., Gilbert, D., Song, H., Guoli, M., Wen, Z., Brinton, M., Gunjan, A., & Tang, H. Zika virus infection induces DNA damage response and S-phase arrest in human neural progenitor cells that enhances viral replication. Journal of Virology, Sep 2019, 93 (20) e00638-19; DOI: 10.1128/JVI.00638-19. (Impact factor: 4.501)},
{Phillips, E. O. N., Giovinazzi, S., Menz, S. L., \textbf{Son, Y.}, Gunjan, A. Preparation of Cell
Extracts by Cryogrinding in an Automated Freezer Mill. Journal of Visualized Experiments (167), 2021, e61164, doi:10.3791/61164. (Impact factor: 1.2)},
{(In preparation} \textbf{Son, Y.} and Gunjan, A. Identification of critical cellular pathways involved in keloid
formation and response to therapy using comprehensive whole-genome sequencing based approaches.},
{(Submitted) Burkhart, S. L., Giovinazzi, S., Phillips, E., Singh, R. K., \textbf{Son, Y.}, Canzani, D., Kabbaj,
Marie-Helene M., Leake, N., Horabin, J., Paik, J., and Gunjan, A. Histone variant H3.3 plays an evolutionarily conserved role in DNA repair that can be targeted for cancer therapy.}

%---------------------------------------------------------------------------------------
%	Patent SECTION
%--------------------------------------------------------------------------------------
\cvsection{Patent applications (US)}{
{Gunjan, A., \textbf{Son, Y.}, & Phillips, E. Methods for predicting steroid responsiveness.
(Application-62/991,423).}
}

%---------------------------------------------------------------------------------------
%	Honors/Activities SECTION
%--------------------------------------------------------------------------------------
\cvsection{Honors and Activities}{
{7 poster presentations for conferences},
{3 oral presentations for conferences},
{Scholarship, Advanced Sequencing Technologies & Bioinformatics Analysis in the Cold Spring Harbor Laboratory (2020)},
{Scholarship, 24th Summer Institute in Statistical Genetics (SISG) in University of Washington in Seattle (2019)},
{Travel Scholarship, International Mammalian Genome Conference (2018)},
{Best Student poster, Southern Genome Maintenance Conference (2018)},
{Scholarship, Short Course On Next-Generation Sequencing in University of Alabama in Birmingham (2017)}
}

%---------------------------------------------------------------------------------------
%	Skills
%----------------------------------------------------------------------------------------
\cvsection{Skills}

\cvevent{Genomics}{
	{NGS library preparation bench works (RNAseq, WGS, WGBS)},
	{WGS analysis to assess and quantify variant calling accuracy for different variant types (SNVs, insertions, deletions, CNVs, SVs, etc.,)},
	{RNAseq analysis to determine differentially expressed genes and identify differential transcript usages},
	{Methlome-seq analysis to find epigenetic markers of diseases and SNPs}
}

%
\cvevent{Molecular and Cellular biology}{
	{RT-qPCR, RNA and DNA extractions, Molecular cloning, Protein and histone extractions, Immunofluorescence microscopy, histology}
}

%\textcolor{softcol}{\hrule}

%
\cvevent{Computer skills}{
	{R, Python, Bash, and Linux for HPC}
}

%\textcolor{softcol}{\hrule}


%--------------------------------------------------------------------------------------------------
%	ARTIFICIAL FOOTER (fancy footer cannot exceed linewidth) 
%--------------------------------------------------------------------------------------------------

\null
\vspace*{\fill}
\hspace{-0.25\linewidth}\colorbox{white}{\makebox[1.5\linewidth][c]{\mystrut  \textnormal{\textcolor{sectcol}{www.jankuester.com} $\cdot$ \textcolor{sectcol}{github.com/jankapunkt}}}}




%============================================================================%
%
%
%
%	DOCUMENT END
%
%
%
%============================================================================%
\end{document}
